\itemhead{2.3 例题 · 例题5（教材p.55）}{XJ-B1-C02-EX-P055-Y0177}
\label{XJ-B1-C02-EX-P055-Y0177}
{\footnotesize 湘教版教材p.55 / PDF第62页 / 2.3 例题 / 例题5}
\field{原题}{若对任意的实数 $x$，一元二次不等式 $x^2+2(1+k)x+3+k>0$ 恒成立，求实数 $k$ 的取值范围。}
\field{参考答案}{由 $\Delta=4(k^2+k-2)<0$ 得 $k\in(-2,1)$，端点不能取。}
{\small \field{关联知识}{I06 恒成立、有解与范围内最值（教材帮p.71、72、73 / PDF80、81、82）}
\field{方法与易错点}{先明确主元及参数，分降次情况，再用判别式、配方或区间最值。 易错：严格号与非严格号混淆；把有解条件当恒成立；直接用全域判别式处理有限区间。}
\field{答案依据}{例题原解位置：湘教版教材p.55及题库原图；独立编写可编辑参考解答。答案册此部分主要收练习、习题，不伪造例题答案册位置。}}
